\section{Introduction}


Large Language Models (LLMs) have significantly transformed the landscape of artificial intelligence and software engineering. While proprietary models such as GPT and Claude dominate the consumer market, public open-source models such as DeepSeek are increasingly gaining popularity among business entities due to their enhanced transparency, lower deployment costs, reduced legal and licensing constraints, and greater flexibility for customization. Indeed, popular model-sharing platforms (e.g., Hugging Face) host over 2 million models~\cite{hg}, with top models receiving more than one million downloads~\cite{Jiang2023AnES}.

Notably, a closer examination of the interconnections among these models reveals a significant degree of reuse~\cite{Jiang2023AnES}. In particular, many models are derived from existing base models through techniques such as fine-tuning~\cite{Fu2022OnTE}, model merging~\cite{Arora2024HeresAF}, and quantization~\cite{Yang2025ASO}. Such derivations enable the adaptation of general-purpose models to domain-specific scenarios, capability enhancement, and efficiency optimization. However, while enjoying the benefits, the model reusing also bring the inherited risks~\cite{chen2026empirical, wang2025large, hu2025large, egashira2025mind}, including legal liabilities (e.g., Intellectual Property violations) to security threats~\cite{jiang2023empirical, jia2022badencoder, greshake2023not} (e.g., backdoors, jailbreaks, and prompt injection).

An effective strategy for risk mitigation is to ensure transparency in model composition, such as using AI Bills of Materials (AIBOM)~\cite{aibom}. However, although platforms provide mechanisms (e.g., model cards in Hugging Face~\cite{hgmodelcard}) that allow contributors to document meta-information, such as model reuse details or model performance~\cite{Mitchell2018ModelCF}, such practices are not consistently adopted proactively by model developers. This leads to incomplete disclosure of model information. For instance, prior work reports that the ninth most downloaded pretrained model on Hugging Face, with over 6.9 million monthly downloads, lacks a model card~\cite{Jiang2023AnES}. Consequently, it undermines model transparency, thereby hindering effective risk traceability and mitigation. 

To address this problem, a range of model fingerprinting techniques can be leveraged. While dynamic fingerprinting approaches rely on analyzing a model's inference behavior in response to carefully designed input queries~\cite{guan2022you, peng2022fingerprinting, yang2022metafinger, pasquini2025llmmap, oyama2025mapping, wu2025llm, gubri2024trap, yamabe2025mergeprint, xu2026dnf, zhang2024reef, wu2025tensorguard, xu2025delta}, they are often sensitive to input selection and may suffer from limited robustness under model modifications such as fine-tuning or compression. Therefore, we focus on the problem of \textit{model lineage construction via static fingerprints}. 

\textbf{Empirical Study.} We begin with an empirical study to understand the dependency prevalence and dependency types of public model reuse.


The problem of opaque model lineage is increasingly pressing. LLMs are often trained on massive datasets collected from multiple sources, combined with pre-trained models, and further fine-tuned in proprietary pipelines. This complexity obscures the relationship between models, making it difficult to trace the origin of errors, detect unintended biases, or comply with licensing requirements. Consequently, software engineers, AI practitioners, and policymakers face challenges in maintaining trust, reliability, and accountability in AI-driven systems. Despite the growing literature on supply chain analysis in traditional software and deep learning frameworks, systematic studies focused on LLM lineage remain sparse.

This paper seeks to address this gap by investigating how far current methods go in tracing and analyzing the lineage of LLMs. Specifically, we ask: (1) What techniques exist to track dependencies and derivations among LLMs? (2) How effective are these techniques in providing visibility into model origins and transformations? (3) What limitations persist, and where is further research needed? By framing LLM lineage in the context of software engineering, we aim to connect advances in AI with established practices in software maintenance, traceability, and quality assurance.
